% =============================================================================
% 00_assumption_appendix.tex
% Assumption Validation Appendix -- Portfolio VaR Project
% Compile: pdflatex -> bibtex -> pdflatex -> pdflatex
% =============================================================================

\documentclass[11pt,a4paper]{article}

\usepackage[utf8]{inputenc}
\usepackage[T1]{fontenc}
\usepackage{lmodern}
\usepackage[margin=2.5cm]{geometry}

% Math
\usepackage{amsmath,amssymb,amsthm,mathtools}

% Tables
\usepackage{booktabs}
\usepackage{longtable}
\usepackage{array}
\usepackage{tabularx}
\usepackage{makecell}

% Code listings
\usepackage{listings}
\usepackage{xcolor}
\definecolor{codegray}{rgb}{0.95,0.95,0.95}
\definecolor{codegreen}{rgb}{0,0.5,0}
\lstdefinestyle{pycode}{
    backgroundcolor=\color{codegray},
    commentstyle=\color{codegreen},
    basicstyle=\ttfamily\footnotesize,
    breakatwhitespace=false,
    breaklines=true,
    captionpos=b,
    keepspaces=true,
    language=Python,
    showspaces=false,
    showstringspaces=false,
    showtabs=false,
    tabsize=2
}
\lstset{style=pycode}

% Hyperref (load last among the package list, before cleveref)
\usepackage[hidelinks]{hyperref}
\usepackage{cleveref}

% Custom commands
\newcommand{\asmhead}[2]{\paragraph{#1 --- #2}\mbox{}\\}
\newcommand{\asmfield}[1]{\medskip\noindent\textbf{#1:}\ }
\newcommand{\code}[1]{\texttt{#1}}
\newcommand{\validated}{\textbf{Validated}}
\newcommand{\diagnostic}{\textbf{Diagnostic check only}}
\newcommand{\notvalidated}{\textbf{NOT validated}}

\title{Assumption Validation Appendix\\[0.5em]
\large Portfolio Value-at-Risk Project\\
\normalsize Frankfurt School of Finance \& Management}
\author{Quantitative Risk Team}
\date{\today}

% ---------------------------------------------------------------------------
\begin{document}
\maketitle
\tableofcontents
\newpage

% ===========================================================================
\section{Executive Summary}\label{sec:executive-summary}
% ===========================================================================

This appendix enumerates every modeling assumption in the portfolio
Value-at-Risk project, classifies each by type, documents what validation
(if any) was performed in-project, and states the concrete consequence if
the assumption fails.

\medskip
\noindent\textbf{Total assumptions documented:} 40, across seven categories.

\medskip
\begin{tabular}{@{}lrr@{}}
\toprule
\textbf{Validation status} & \textbf{Count} & \textbf{\%} \\
\midrule
Formally validated (test with $p$-value or algebraic proof) & 9  & 22.5 \\
Diagnostic check only (plot / qualitative inspection)       & 16 & 40.0 \\
NOT validated (accepted without testing)                    & 15 & 37.5 \\
\bottomrule
\end{tabular}

\medskip
\noindent\textbf{Three assumptions whose failure would most affect reported VaR:}

\begin{enumerate}
  \item \textbf{\cref{asm:C10} --- Returns i.i.d.\ (no GARCH filtering in
    copula model).} Raw log-returns exhibit well-documented volatility
    clustering. Without GARCH pre-filtering, pseudo-observations are \emph{not}
    i.i.d., which distorts the copula correlation estimate and the marginal
    fits. During high-volatility regimes (GFC~2008, COVID~2020) this can
    cause 30--50\% underestimation of VaR. This is explicitly flagged as a
    known limitation of the course methodology.

  \item \textbf{\cref{asm:A4} --- Losses i.i.d.\ within 500-day rolling
    window.} Both the plain EVT and the copula model assume that the 500
    daily loss observations used to estimate tail parameters are i.i.d.
    Volatility clustering, mean-reversion, and regime changes all violate
    this. The EVT backtest produced 56 exceptions against 42.7 expected
    ($p_{\mathrm{IND}} = 0.0064$), confirming temporal clustering of
    breaches.

  \item \textbf{\cref{asm:D5} --- Historical sample period representative
    of future regimes.} The 2006--2024 sample spans GFC, COVID, and the
    2022 rate-hike cycle but excludes hyperinflationary episodes, sovereign
    default cascades, and USD-reserve crises. Any VaR estimate is
    conditional on this historical distribution; a structurally different
    future would invalidate every model simultaneously.
\end{enumerate}

% ===========================================================================
\section{Methodology}\label{sec:methodology}
% ===========================================================================

\subsection{Scope discovery}\label{sec:scope-discovery}

Assumptions were enumerated by systematic line-by-line reading of the
following source files (not from memory):

\begin{itemize}
  \item \code{src/var\_methods/evt.py} --- EVT/POT VaR (508~lines)
  \item \code{src/var\_methods/garch\_evt.py} --- GARCH+EVT VaR (705~lines)
  \item \code{src/var\_methods/var\_copula.py} --- copula MC VaR ($\approx$880~lines)
  \item \code{backtesting/backtest.py} --- Kupiec + Christoffersen framework (433~lines)
  \item \code{src/data/compute\_returns.py} --- log-return computation
  \item \code{src/data/compute\_pnl.py} --- IRS + straddle P\&L
  \item \code{src/data/portfolio\_pricing.py} --- Black-Scholes + IRS pricing
  \item \code{src/data/config.py} --- portfolio constants
\end{itemize}

Files \emph{not} scanned (with justification):
\begin{itemize}
  \item \code{src/var\_methods/delta\_normal.py}, \code{historical\_simulation.py} ---
    belong to other team members; not part of this appendix's scope.
  \item \code{src/data/download\_data.py} --- data download utility; no
    modeling assumptions beyond data alignment (covered under D.6).
  \item Plotting helpers (\code{plot\_*} functions within each VaR module) ---
    purely presentational; no statistical decisions embedded.
  \item \code{docs/pricing/\_generate\_docx.py} --- documentation generator; no model logic.
\end{itemize}

\noindent Ten or more modeling components were identified, exceeding the
five-component minimum threshold. Discovery is complete.

\subsection{Category taxonomy}\label{sec:taxonomy}

Every assumption belongs to exactly one of the following categories:

\medskip
\begin{tabular}{@{}p{4cm}p{10.5cm}@{}}
\toprule
\textbf{Category} & \textbf{Examples} \\
\midrule
A.\ Distributional & Returns GBM / log-normal; Student-$t$ innovations;
  GPD above threshold $u$; $t$-copula $\nu=4$ \\[4pt]
B.\ Independence / temporal & Returns i.i.d.; standardised residuals i.i.d.;
  pseudo-observations independent \\[4pt]
C.\ Stationarity & Mean/variance constant in window; stable correlation;
  stable tail shape $\xi$ \\[4pt]
D.\ Parameter / model & GARCH(1,1) sufficient; GPD threshold $= 90$th pct;
  copula df fixed; equal weights \\[4pt]
E.\ Market / economic & No transaction costs; continuous trading; constant
  $r_f$; flat term structure; no jumps \\[4pt]
F.\ Numerical / implementation & MC error negligible at $N{=}20{,}000$;
  GARCH MLE converges; sample $R \approx$ MLE \\[4pt]
G.\ Scope / data & \code{pnl\_total} covers all risks; nonlinear
  independent of linear; sample representative \\
\bottomrule
\end{tabular}

\subsection{Citation conventions}\label{sec:citations}

References to Irle's lecture notes follow the page numbering of the
2025 edition distributed at Frankfurt School. All citations use
BibTeX keys listed in \code{references.bib}. Backtest statistics
quoted in assumption entries are taken directly from
\code{outputs/tables/backtest\_evt.csv},
\code{outputs/tables/backtest\_garch\_evt.csv}, and
\code{outputs/tables/backtest\_copula.csv}.

% ===========================================================================
\section{Per-Assumption Entries}\label{sec:entries}
% ===========================================================================

% ---------------------------------------------------------------------------
\subsection{Category A --- Distributional}\label{sec:catA}
% ---------------------------------------------------------------------------

\subsubsection{A.1 --- Pickands-Balkema-de Haan convergence}\label{asm:A1}

\asmfield{Statement}
The conditional excess distribution of losses $L$ above a sufficiently
high threshold $u$ converges to the Generalised Pareto Distribution
$G_{\xi,\sigma}(x) = 1 - (1 + \xi x/\sigma)^{-1/\xi}$ as $u \to \infty$.

\asmfield{Source in code}
\code{evt.py:220-272} and \code{garch\_evt.py:391-445} ($\hat{Z}_t$ residuals):
\begin{lstlisting}
u = np.quantile(losses_w, threshold_q)
exceedances = losses_w[losses_w > u] - u
c, loc, sigma = genpareto.fit(exceedances, floc=0)
\end{lstlisting}

\asmfield{Theoretical basis}
\cite{Pickands1975}; \cite{BalkemadeHaan1974}. Irle lecture notes,
p.~213 (Theorem~9.1). \cite{McNeilFreyEmbrechts2005}, \S~7.2.

\asmfield{Validation in this project}
\diagnostic\ --- Mean Excess Plot (MEP) and Hill estimator plot generated
in \code{plot\_threshold\_diagnostics()} (\code{evt.py:107-183}). A linear
MEP region above the chosen threshold supports GPD convergence
qualitatively. KS goodness-of-fit computed per rolling window
(\code{evt.py:270}); zero windows with KS $p < 0.05$ recorded in the
4{,}269-day production run. No formal convergence test was performed
because no such test exists with finite samples.

\asmfield{Consequence if assumption fails}
VaR is \emph{underestimated} for true tail behaviour heavier than GPD
(e.g., jump-diffusion or $\alpha$-stable processes). Approximate
magnitude: 10--30\% underestimation at the 99.9\% confidence level based
on empirical POT failures documented in \cite{McNeilFreyEmbrechts2005},
\S~7.3. At the 99\% level used in this project, the impact is smaller
($\approx 5$--$15\%$) because 99\% is closer to the body of the
distribution.

\asmfield{Mitigation in this project}
Regulatory floor $\max(\text{VaR}, u, 0)$ imposed (\code{evt.py:262-264});
$\xi$ clamped to $[-0.5, 1.0]$ (\code{evt.py:238-249}) to prevent
infinite-mean fits. Neither mitigation restores statistical validity if GPD
convergence fails --- both are conservative policy choices.

% - - - - - - - - - - - - - - - - - - - - - - - - - - - - - - - - - - - - -
\subsubsection{A.2 --- 90th percentile is a valid GPD threshold}\label{asm:A2}

\asmfield{Statement}
The 90th percentile of the rolling 500-day loss distribution is a
sufficiently high threshold $u$ such that Assumption~\cref{asm:A1}
holds, yielding approximately $N_u = 50$ exceedances per window.

\asmfield{Source in code}
\code{evt.py:70}: \lstinline|THRESHOLD_Q = 0.90|;
\code{garch\_evt.py:94}: \lstinline|THRESHOLD_Q = 0.90|;
\code{var\_copula.py} does not use POT directly.

\asmfield{Theoretical basis}
Irle, p.~215 (MEP linearity criterion). \cite{McNeilFreyEmbrechts2005},
\S~7.2.2 recommends $N_u \ge 50$ exceedances for stable GPD MLE.
With 500 observations and $q = 0.90$, $N_u \approx 50$ exactly.

\asmfield{Validation in this project}
\diagnostic\ --- Full-sample MEP at \code{outputs/figures/evt\_01\_threshold\_diagnostics.png}
shows a roughly linear region beginning near the 85th--90th percentile,
consistent with GPD applicability. No formal threshold selection test
(e.g., automated stopping rule) was applied.

\asmfield{Consequence if assumption fails}
\emph{Threshold too low} (below GPD convergence zone): exceedances are not
GPD-distributed; $\hat{\xi}$ is biased downward and VaR is
\emph{underestimated} by $5$--$15\%$. \emph{Threshold too high}
($q > 0.95$, few exceedances): high estimation variance; $\hat{\xi}$ is
noisy; VaR fluctuates excessively. The 90th percentile is the standard
industry choice for 500-day windows.

\asmfield{Mitigation in this project}
\code{MIN\_EXCEEDANCES = 10} guard in \code{evt.py:71} ensures empirical
fallback if threshold is so high that fewer than 10 exceedances remain.
Fallback triggered zero times in the production run.

% - - - - - - - - - - - - - - - - - - - - - - - - - - - - - - - - - - - - -
\subsubsection{A.3 --- $\xi \in [-0.5, 1.0]$ is a plausible range}\label{asm:A3}

\asmfield{Statement}
The GPD shape parameter $\xi$ for financial daily loss tails lies in
$[-0.5, 1.0]$: $\xi > 1$ implies infinite mean (pathological);
$\xi < -0.5$ implies implausibly bounded losses.

\asmfield{Source in code}
\code{evt.py:238-249}:
\begin{lstlisting}
if xi > 1.0:
    warnings.warn(...)
    xi = 1.0
elif xi < -0.5:
    warnings.warn(...)
    xi = -0.5
\end{lstlisting}
Identical logic in \code{garch\_evt.py:407-418}.

\asmfield{Theoretical basis}
Irle, p.~212. Cont \& Tankov (2004), p.~93. \cite{McNeilFreyEmbrechts2005},
\S~7.2 notes $\xi \in (0, 0.5)$ is typical for equity daily returns.
Mean rolling $\xi = 0.0475$ (EVT) and $\xi = 0.0312$ (GARCH+EVT) are
both well within bounds.

\asmfield{Validation in this project}
\diagnostic\ --- Rolling $\xi$ plotted in \code{outputs/figures/evt\_02\_var\_and\_xi.png}.
Clamping warnings emitted when triggered (\code{warnings.warn}); no
clamping events recorded in the 4{,}269-day run (mean $\xi \ll 0.5$,
never near either bound). Literature support is strong but no formal
hypothesis test was applied.

\asmfield{Consequence if assumption fails}
If true $\xi > 1$ (not clamped): VaR formula involves $\sigma/\xi$, which
could produce VaR estimates orders of magnitude too large. Clamping at 1.0
\emph{overestimates} VaR in this case (conservative). If true $\xi < -0.5$:
bounded tail is real but we treat it as less extreme --- VaR is
\emph{overestimated} (also conservative). Both failures produce \emph{safe}
over-estimation.

\asmfield{Mitigation in this project}
Clamping is itself the mitigation. Audit trail via \code{warnings.warn}.

% - - - - - - - - - - - - - - - - - - - - - - - - - - - - - - - - - - - - -
\subsubsection{A.4 --- Losses are i.i.d.\ within a 500-day rolling window}\label{asm:A4}

\asmfield{Statement}
The 500 daily portfolio losses used to fit the GPD (EVT) or the marginal
distributions (copula) are independent and identically distributed, i.e.,
no serial dependence, no regime change, and constant parameters within
each window.

\asmfield{Source in code}
\code{evt.py:309-322} (rolling loop assumes i.i.d.\ losses);
\code{var\_copula.py:387-388} (pseudo-observations assume i.i.d.\ returns).

\asmfield{Theoretical basis}
Irle, p.~213 (Theorem~9.1 states the i.i.d.\ condition explicitly).
\cite{McNeilFrey2000} motivate GARCH pre-filtering precisely because this
assumption fails for raw returns.

\asmfield{Validation in this project}
\notvalidated\ --- No Ljung-Box, ARCH-LM, or runs test was performed on
the rolling windows. The EVT backtest confirms failure: 56 exceptions
($1.31\%$) vs.\ 42.7 expected ($1.00\%$); Christoffersen LR$_{\mathrm{IND}}
= 7.42$, $p = 0.0064$ --- independent exceptions rejected at $5\%$
(\code{outputs/tables/backtest\_evt.csv}).

\asmfield{Consequence if assumption fails}
Exception clustering causes VaR to be \emph{underestimated during crises}
and \emph{overestimated during calm periods}. Magnitude: the EVT exception
rate during GFC 2008 exceeded $1.31\%$ annually --- consistent with
$20$--$50\%$ VaR underestimation during high-volatility regimes. The
rolling window partially mitigates by discarding data more than 500 days
old, but cannot correct for intra-window clustering.

\asmfield{Mitigation in this project}
GARCH+EVT eliminates this assumption for GARCH-filtered residuals (see
\cref{asm:B4}). For plain EVT and the copula, rolling window reduces
staleness but does not restore i.i.d.\ --- flagged as known limitation.

% - - - - - - - - - - - - - - - - - - - - - - - - - - - - - - - - - - - - -
\subsubsection{A.5 --- KS $p$-value informative despite anti-conservative bias}\label{asm:A5}

\asmfield{Statement}
The KS goodness-of-fit test applied to GPD exceedances, with parameters
estimated from the same data, produces $p$-values that are informative as
relative quality indicators across rolling windows even though they are
biased upward (anti-conservative).

\asmfield{Source in code}
\code{evt.py:266-270}:
\begin{lstlisting}
# parameters are estimated from the same data -- anti-conservative
_, ks_pvalue = kstest(exceedances, "genpareto", args=(xi, 0, sigma))
\end{lstlisting}
Identical caution in \code{garch\_evt.py:433-434}.

\asmfield{Theoretical basis}
\cite{McNeilFreyEmbrechts2005}, \S~7.3.1 note that KS with estimated
parameters has inflated $p$-values. A corrected test (e.g., Anderson-Darling
with bootstrapped critical values) would be more rigorous.

\asmfield{Validation in this project}
\notvalidated\ --- The bias is documented in code comments but not corrected.
No poor-fit windows (KS $p < 0.05$) were recorded in production, which
likely reflects the anti-conservative bias rather than true perfect fit.

\asmfield{Consequence if assumption fails}
KS $p$-values are systematically too high; a genuinely poor GPD fit may
pass the test. True rejection rate could be 20--40\% higher than reported.
Used only as a relative quality indicator (\code{ks\_pvalue} column in
output CSVs), not as an acceptance gate.

\asmfield{Mitigation in this project}
None within the current implementation. The code comment documents the
bias explicitly (\code{evt.py:267-269}).

% - - - - - - - - - - - - - - - - - - - - - - - - - - - - - - - - - - - - -
\subsubsection{A.6 --- Empirical quantile is an acceptable fallback}\label{asm:A6}

\asmfield{Statement}
When GPD MLE fails or fewer than \code{MIN\_EXCEEDANCES = 10} exceedances
are available, the empirical 99th percentile of the loss window is an
acceptable VaR estimate.

\asmfield{Source in code}
\code{evt.py:224-226}:
\begin{lstlisting}
if N_u < MIN_EXCEEDANCES:
    fb = np.quantile(losses_w, alpha)
    return fb, np.nan, np.nan, fb
\end{lstlisting}

\asmfield{Theoretical basis}
No theoretical justification --- pragmatic fallback. The empirical quantile
is a consistent estimator of the population quantile but cannot extrapolate
beyond the maximum observed loss.

\asmfield{Validation in this project}
\validated\ by absence of activation: 0 fallback days recorded out of 4{,}269
(\code{outputs/run log}).

\asmfield{Consequence if assumption fails}
VaR cannot exceed the maximum observed loss in the window --- a severe
limitation if the true 99th percentile lies beyond the sample maximum.
Magnitude: unknown, but potentially large for very short tails.

\asmfield{Mitigation in this project}
The threshold choice ($q = 0.90$) ensures $N_u \approx 50 \gg 10$,
making the fallback essentially never triggered in production.

% - - - - - - - - - - - - - - - - - - - - - - - - - - - - - - - - - - - - -
\subsubsection{A.7 --- Regulatory floor is a defensible policy choice}\label{asm:A7}

\asmfield{Statement}
$\text{VaR}_{\text{floored}} = \max(\text{VaR}_{\text{raw}}, u, 0)$ is a
conservative regulatory-reporting policy choice, \emph{not} a statistical
claim about the true tail quantile.

\asmfield{Source in code}
\code{evt.py:262-264}:
\begin{lstlisting}
var_floored = max(var_raw, u)    # conservative: never below threshold
var_floored = max(var_floored, 0.0)  # defensive: non-negative
\end{lstlisting}

\asmfield{Theoretical basis}
Irle (regulatory context). When $\xi < 0$ the GPD has a bounded support
and the true VaR can legitimately fall below $u$. The floor is a
deliberate policy override, not a correction of a statistical error.

\asmfield{Validation in this project}
\validated\ --- Separation of \code{VaR\_EVT\_raw} and \code{VaR\_EVT}
(floored) implemented per model validation commit A7/B9. Zero days where
the floor was binding in the production run (mean floor impact: \$0).

\asmfield{Consequence if assumption fails}
If the floor is too aggressive and the true VaR is below $u$: VaR is
overestimated (conservative --- acceptable for regulatory reporting). If
the floor is removed and $\xi < 0$ produces VaR $< 0$ (numerical artefact):
VaR becomes meaningless. The \code{var\_raw} diagnostic column permits
post-hoc analysis without the floor.

\asmfield{Mitigation in this project}
Both \code{VaR\_EVT} (floored) and \code{VaR\_EVT\_raw} (unfloored)
stored in \code{data/processed/var\_evt.csv}.

% ---------------------------------------------------------------------------
\subsection{Category B --- Independence and Temporal Structure}\label{sec:catB}
% ---------------------------------------------------------------------------

\subsubsection{B.1 --- GARCH(1,1) is sufficient}\label{asm:B1}

\asmfield{Statement}
A GARCH(1,1) model --- one lag of past squared returns and one lag of
past variance --- adequately captures the conditional heteroskedasticity
in the portfolio return series.

\asmfield{Source in code}
\code{garch\_evt.py:212-213}:
\begin{lstlisting}
garch = arch_model(train, vol="Garch", p=1, q=1,
                   dist="t", mean="Constant")
\end{lstlisting}

\asmfield{Theoretical basis}
\cite{Bollerslev1986}. Irle, Section~8.3 (p.~172-186). Extensive
empirical literature (e.g., Hansen \& Lunde 2005) confirms GARCH(1,1)
dominates more complex specifications for daily financial returns.

\asmfield{Validation in this project}
\diagnostic\ --- Standardised residual ACF and QQ-plot generated in
\code{outputs/figures/garch\_evt\_01\_diagnostics.png}. GARCH+EVT
backtest: 41 exceptions vs.\ 42.7 expected, Kupiec $p = 0.7936$, CC
$p = 0.6914$ (\code{outputs/tables/backtest\_garch\_evt.csv}) --- best
backtest result of the three methods. No AIC/BIC comparison vs.\
GARCH(2,1) performed.

\asmfield{Consequence if assumption fails}
Direction is ambiguous. GARCH(2,1) could reduce residual autocorrelation,
potentially improving both VaR accuracy and Christoffersen independence
test. Magnitude: $< 10\%$ impact on VaR based on typical GARCH order
mis-specification studies.

\asmfield{Mitigation in this project}
None --- sensitivity test (see \cref{sec:sensitivity-roadmap}) would
quantify the impact.

% - - - - - - - - - - - - - - - - - - - - - - - - - - - - - - - - - - - - -
\subsubsection{B.2 --- GARCH innovations are Student-$t$}\label{asm:B2}

\asmfield{Statement}
Conditional on $\sigma_t$, portfolio returns $X_t = \sigma_t Z_t$ where
$Z_t \sim t(\nu)$ (Student-$t$ with $\nu$ estimated jointly with GARCH
parameters).

\asmfield{Source in code}
\code{garch\_evt.py:213}: \lstinline|dist="t"| in \code{arch\_model};
$\nu$ estimated in \code{res.params}.

\asmfield{Theoretical basis}
Irle, p.~179-180. Fat tails in daily returns are well-documented; Normal
GARCH underestimates tail probabilities; Student-$t$ GARCH is standard
industry practice.

\asmfield{Validation in this project}
\diagnostic\ --- QQ-plot of standardised residuals vs.\ Student-$t$ in
\code{outputs/figures/garch\_evt\_01\_diagnostics.png}. Residual tails
deviate slightly from Student-$t$ (justifying the subsequent EVT step).

\asmfield{Consequence if assumption fails}
If true innovations are heavier-tailed than Student-$t$ (e.g., stable
$\alpha < 2$): GPD step absorbs residual tail risk, so impact on final
VaR is small ($< 5\%$ at 99\%). Normal innovations would miss more.

\asmfield{Mitigation in this project}
EVT on standardised residuals in Step~B provides a non-parametric tail
model that does not rely on the Student-$t$ assumption (see \cref{asm:B8}).

% - - - - - - - - - - - - - - - - - - - - - - - - - - - - - - - - - - - - -
\subsubsection{B.3 --- GARCH stationarity ($\alpha + \beta < 1$)}\label{asm:B3}

\asmfield{Statement}
The GARCH(1,1) parameters satisfy $\alpha_1 + \beta < 1$ (covariance
stationarity; Irle p.~176), ensuring finite unconditional variance.

\asmfield{Source in code}
\code{garch\_evt.py:239-245}:
\begin{lstlisting}
ab = alpha1_d + beta_d
flag = "OK" if ab < 1.0 else "!! UNIT ROOT"
print(f"  alpha_1 + beta = {ab:.6f}  -> {flag}")
\end{lstlisting}

\asmfield{Theoretical basis}
Irle, p.~176. \cite{Bollerslev1986}. Stationarity is required for
consistent GARCH MLE and for the conditional VaR formula
$\text{VaR}_t = V_0 \cdot \hat{\sigma}_t \cdot q_{\mathrm{EVT}}(\alpha)$
to be well-defined.

\asmfield{Validation in this project}
\validated\ --- Initial fit (at $t = 500$): $\alpha_1 + \beta = 0.9718$.
Across 86 expanding-window refits: range 0.971--0.997, all $< 1$.
Output log confirms \code{"OK"} for every refit.

\asmfield{Consequence if assumption fails}
$\alpha_1 + \beta \ge 1$: integrated GARCH (IGARCH); conditional variance
is non-stationary; $\hat{\sigma}_t \to \infty$ over long horizons; VaR
diverges. In the backtest period this did not occur.

\asmfield{Mitigation in this project}
\code{arch} library enforces stationarity constraints during MLE by
default. Diagnostic flag printed for every refit.

% - - - - - - - - - - - - - - - - - - - - - - - - - - - - - - - - - - - - -
\subsubsection{B.4 --- Standardised residuals $\hat{Z}_t$ are approximately i.i.d.}\label{asm:B4}

\asmfield{Statement}
After GARCH filtering, $\hat{Z}_t = (X_t - \hat{\mu})/\hat{\sigma}_t$
are approximately independently and identically distributed, justifying
application of POT/GPD to the lower tail of $\{\hat{Z}_t\}$.

\asmfield{Source in code}
\code{garch\_evt.py:267-268}:
\begin{lstlisting}
z_arr[t-1] = (r_vals[t-1] - mu_hat_prev) / cond_vol_arr[t-1]
\end{lstlisting}
The \code{mu\_hat\_prev} fix (commit B10) ensures residuals are computed
under the same parameter regime as their conditional volatility, removing
a cross-regime contamination artefact at refit boundaries.

\asmfield{Theoretical basis}
\cite{McNeilFrey2000}, \S~2. Irle, Section~8.3. The two-step approach is
motivated precisely by the approximate i.i.d.\ property of GARCH residuals.

\asmfield{Validation in this project}
\diagnostic\ --- QQ-plot of $\hat{Z}_t$ vs.\ Student-$t$ at
\code{outputs/figures/garch\_evt\_01\_diagnostics.png}. Residual mean
$= -0.0176$ (expected $\approx 0$), standard deviation $= 0.9936$
(expected $\approx 1$). No formal Ljung-Box test on $\hat{Z}_t$ performed.

\asmfield{Consequence if assumption fails}
Residual autocorrelation: GPD fit on clustered $\hat{Z}_t$ produces a
biased $\hat{\xi}$; VaR is underestimated during volatility bursts. The
excellent GARCH+EVT backtest (Kupiec $p = 0.7936$, CC $p = 0.6914$)
suggests the approximation is adequate.

\asmfield{Mitigation in this project}
GARCH filtering is the primary mitigation (\cref{asm:B8}). B10 fix
eliminates a secondary source of residual bias at refit boundaries.

% - - - - - - - - - - - - - - - - - - - - - - - - - - - - - - - - - - - - -
\subsubsection{B.5 --- GARCH refit every 50 days is sufficient}\label{asm:B5}

\asmfield{Statement}
Re-estimating GARCH parameters on an expanding window every 50 trading
days captures time-variation in parameters adequately; parameters are
stable within each 50-day inter-refit interval.

\asmfield{Source in code}
\code{garch\_evt.py:97}: \lstinline|REFIT_EVERY = 50|;
\code{garch\_evt.py:202}: \lstinline|(t - window) % refit_every == 0|.

\asmfield{Theoretical basis}
No formal theoretical guidance. Irle, Section~8.3 discusses rolling vs.\
expanding estimation without recommending a specific cadence.

\asmfield{Validation in this project}
\notvalidated\ --- No sensitivity test comparing REFIT\_EVERY $\in
\{25, 100, 250\}$ was performed. Parameter time series show $\alpha_1 +
\beta$ stable in 0.971--0.997 across 86 refits, suggesting moderate
stability, but inter-refit drift is not formally quantified.

\asmfield{Consequence if assumption fails}
Fast regime change (e.g., COVID onset) within a 50-day window: GARCH
parameters may be stale by 25--50 days; $\hat{\sigma}_t$ under- or
overestimates true volatility; VaR biased accordingly. Magnitude $< 5\%$
based on observed parameter stability across refits.

\asmfield{Mitigation in this project}
Expanding window (not rolling) means each refit uses all available history;
parameters change slowly. Sensitivity test listed in
\cref{sec:sensitivity-roadmap}.

% - - - - - - - - - - - - - - - - - - - - - - - - - - - - - - - - - - - - -
\subsubsection{B.6 --- GARCH-fitted $\hat{\mu}$ is more accurate than sample mean}\label{asm:B6}

\asmfield{Statement}
Using the GARCH-MLE conditional mean estimate $\hat{\mu}$ (instead of the
sample mean $\bar{r}$) for computing standardised residuals reduces bias
in $\hat{Z}_t$.

\asmfield{Source in code}
\code{garch\_evt.py:219}: \lstinline|mu_hat = res.params["mu"] / 100.0|.

\asmfield{Theoretical basis}
Standard MLE argument: the GARCH MLE provides a consistent, asymptotically
efficient estimator of $\mu$ jointly with $\omega$, $\alpha_1$, $\beta$,
and $\nu$. Using the sample mean ignores information in the variance
structure.

\asmfield{Validation in this project}
\notvalidated\ --- No comparison between sample-mean and MLE-mean residuals
was performed. In practice, $\hat{\mu} \approx \bar{r} \approx 0$ for daily
financial returns (mean return per day $\approx 0.01\%$); the impact is
negligible.

\asmfield{Consequence if assumption fails}
$|\hat{\mu} - \bar{r}|$ is typically $< 0.005\%$ per day; impact on $\hat{Z}_t$
and hence VaR is $< 0.1\%$.

\asmfield{Mitigation in this project}
GARCH MLE is the standard approach (commit B2).

% - - - - - - - - - - - - - - - - - - - - - - - - - - - - - - - - - - - - -
\subsubsection{B.7 --- $\hat{\sigma}_t$ is the correct conditioning variable}\label{asm:B7}

\asmfield{Statement}
The one-step-ahead GARCH conditional volatility forecast $\hat{\sigma}_t$
(computed from data through $t-1$) is the appropriate scaling variable for
$\text{VaR}_t = V_0 \cdot \hat{\sigma}_t \cdot q_{\mathrm{EVT}}(\alpha)$.

\asmfield{Source in code}
\code{garch\_evt.py:225-227}:
\begin{lstlisting}
fcast = res_t.forecast(horizon=1, reindex=False)
sigma_sq_curr   = fcast.variance.iloc[-1, 0] / 10000.0
cond_vol_arr[t] = np.sqrt(max(sigma_sq_curr, 1e-10))
\end{lstlisting}

\asmfield{Theoretical basis}
\cite{McNeilFrey2000}, Equation~(4): $\text{VaR}_{t|t-1}(\alpha) =
\sigma_t q_Z(\alpha)$. Irle, p.~186.

\asmfield{Validation in this project}
\validated\ --- Structural: the one-step-ahead forecast is the theoretically
correct conditioning variable for next-day VaR by construction. Naming bug
fixed in commit B3 (\code{sigma\_hat\_t1} $\to$ \code{sigma\_hat\_t}) to
reflect this correctly.

\asmfield{Consequence if assumption fails}
If GARCH spec is wrong (see \cref{asm:B1}), $\hat{\sigma}_t$ will also
be wrong; consequences follow from B.1.

\asmfield{Mitigation in this project}
Correctness validated by code review; no separate empirical test needed.

% - - - - - - - - - - - - - - - - - - - - - - - - - - - - - - - - - - - - -
\subsubsection{B.8 --- EVT on residuals reduces volatility-clustering bias}\label{asm:B8}

\asmfield{Statement}
Applying POT/GPD to standardised GARCH residuals $\hat{Z}_t$ rather than
to raw returns $X_t$ removes the volatility-clustering bias in the tail
quantile estimate.

\asmfield{Source in code}
\code{garch\_evt.py:364-445}: \code{\_pot\_var\_residuals()} applied to
\code{z\_w}, not to raw returns.

\asmfield{Theoretical basis}
\cite{McNeilFrey2000}, \S~2--3: the motivation for the two-step procedure
is exactly that GARCH filtering makes residuals approximately i.i.d., after
which EVT can be validly applied.

\asmfield{Validation in this project}
\validated\ by comparison: EVT (plain, 56 exceptions, $p_{\mathrm{IND}} =
0.0064$) vs.\ GARCH+EVT (41 exceptions, $p_{\mathrm{IND}} = 0.4132$).
GARCH pre-filtering eliminates the independence rejection, confirming that
clustering was the driver of EVT's failure.

\asmfield{Consequence if assumption fails}
If GARCH residuals are not approximately i.i.d.\ (see \cref{asm:B4}),
applying EVT to them provides false comfort --- the clustering bias
persists in a more subtle form. The good Christoffersen result suggests
this is not a material concern in this dataset.

\asmfield{Mitigation in this project}
GARCH filtering is the primary mitigation; QQ-plot diagnostic provides
visual confirmation.

% ---------------------------------------------------------------------------
\subsection{Category C --- Stationarity (Copula Model)}\label{sec:catC}
% ---------------------------------------------------------------------------

\subsubsection{C.1 --- Marginals well-fit by the AIC candidate set}\label{asm:C1}

\asmfield{Statement}
For each linear asset (EURUSD, GLD, IEF, SPY), the best-fitting member of
$\{\text{Normal, Logistic, Student-}t\text{, Laplace, GenNormal,
GenLogistic}\}$ under AIC is an adequate model for the return distribution.

\asmfield{Source in code}
\code{var\_copula.py:109-116}: candidate list; \code{var\_copula.py:175-226}:
\code{fit\_best\_marginal()} computes AIC for each candidate.

\asmfield{Theoretical basis}
AIC model selection: Akaike (1974). Irle, \S~6. Student-$t$ marginals are
most commonly selected; mean AIC advantage of best vs.\ Normal: 65.2 units
(\code{outputs/tables/marginal\_selection.csv}).

\asmfield{Validation in this project}
\diagnostic\ --- AIC comparison per refit per asset; most-common winner is
Student-$t$. QQ-plot in \code{outputs/figures/var\_copula\_marginals\_*.png}.
KS $p$-value computed per marginal per refit (\code{var\_copula.py:212-219}).

\asmfield{Consequence if assumption fails}
Poor marginal fit distorts pseudo-observations $U = F(X)$; copula
correlation $R$ is estimated on mis-specified pseudo-observations; simulated
$r_j$ are not from the true marginal. Impact on VaR: $5$--$10\%$ at 99\%
if the true marginal is heavier-tailed than all six candidates.

\asmfield{Mitigation in this project}
Six candidates covering a wide range of tail thicknesses; AIC penalises
over-parameterisation; KS diagnostic stored for review.

% - - - - - - - - - - - - - - - - - - - - - - - - - - - - - - - - - - - - -
\subsubsection{C.2 --- Pseudo-observations are approximately uniform and independent}\label{asm:C2}

\asmfield{Statement}
The rank-based pseudo-observations $U_{i,j} = \text{rank}(X_{i,j}) / (n+1)$
are approximately $\text{Uniform}(0,1)$ and independent across time
(i.e., the marginal dependence structure is correctly removed).

\asmfield{Source in code}
\code{var\_copula.py:231-242}: \code{pseudo\_observations()} with Hazen
correction $U_{i,j} = \text{rank}_j[i] / (n+1)$.

\asmfield{Theoretical basis}
Joe (1997), Ch.~10. \cite{McNeilFreyEmbrechts2005}, \S~5.4.1. Hazen
correction reduces boundary bias vs.\ exact empirical CDF.

\asmfield{Validation in this project}
\diagnostic\ --- Pseudo-observation scatter plots in
\code{outputs/figures/var\_copula\_dependence\_*.png}. Uniformity is
approximately satisfied by construction (rank-based). Serial independence
is \emph{not} validated; see \cref{asm:C10}.

\asmfield{Consequence if assumption fails}
Non-uniform pseudo-observations bias the copula correlation estimate $R$.
Magnitude: $< 5\%$ with correct marginal models. Serial dependence causes
the copula to underestimate tail co-movements during crises.

\asmfield{Mitigation in this project}
Hazen correction mitigates boundary effects; full independence requires
GARCH pre-filtering (not applied in copula model).

% - - - - - - - - - - - - - - - - - - - - - - - - - - - - - - - - - - - - -
\subsubsection{C.3 --- Student-$t$ copula with $\nu = 4$ captures joint dependence}\label{asm:C3}

\asmfield{Statement}
A multivariate Student-$t$ copula with fixed degrees of freedom $\nu = 4$
and correlation matrix $R$ (fitted on pseudo-observations) adequately
captures the joint dependence structure of the four linear assets.

\asmfield{Source in code}
\code{var\_copula.py:93}: \lstinline|NU_T_COPULA = 4|;
\code{var\_copula.py:253-272}: \code{fit\_t\_copula()}.
The code comment at line~29 notes: \emph{``Fitting $\nu$ would be more
rigorous (McNeil-Frey-Embrechts 2005) but departs from the course
methodology.''}

\asmfield{Theoretical basis}
\cite{McNeilFreyEmbrechts2005}, Ch.~7. $t$-copulas with small $\nu$ capture
upper/lower tail dependence $\lambda_U > 0$; Gaussian copulas have
$\lambda_U = 0$.

\asmfield{Validation in this project}
\diagnostic\ (partial) --- CvM GoF test (\cref{asm:C5}) run at every third
refit. CvM rejection rate: 4/18 = 22.2\% (\code{outputs/tables/copula\_gof.csv}).
Three refits (including 2007-12-31 and 2016-12-20) show rejection at $5\%$,
suggesting the $\nu = 4$ model does \emph{not} fit all regimes.

\asmfield{Consequence if assumption fails}
If true $\nu > 4$ (lighter tail dependence than assumed): VaR is
\emph{overestimated} (conservative). If true $\nu < 4$ (heavier tail
dependence): VaR is \emph{underestimated}, particularly for joint tail
events. The 22.2\% rejection rate suggests the fixed $\nu = 4$ is
imperfect. Magnitude: $10$--$20\%$ at 99\% VaR level.

\asmfield{Mitigation in this project}
No corrective action. CvM rejection is documented in output CSVs. Sensitivity
analysis (varying $\nu$) listed in \cref{sec:sensitivity-roadmap}.

% - - - - - - - - - - - - - - - - - - - - - - - - - - - - - - - - - - - - -
\subsubsection{C.4 --- Sample correlation approximates MLE of $R$ for $d \le 4$}\label{asm:C4}

\asmfield{Statement}
Computing the correlation matrix $R$ as the sample correlation of
$t_\nu$-transformed pseudo-observations is a near-optimal approximation
to the true MLE of $R$ for $d \le 4$ assets.

\asmfield{Source in code}
\code{var\_copula.py:253-272}: \code{fit\_t\_copula()}, specifically
\lstinline|R = np.corrcoef(T_emp.T)|.

\asmfield{Theoretical basis}
\cite{McNeilFreyEmbrechts2005}, \S~7.3.3: for $d \le 4$ and fixed $\nu$,
the sample correlation of $t^{-1}_\nu(U)$ is near-optimal. The approximation
error increases with $d$.

\asmfield{Validation in this project}
\notvalidated\ --- No direct comparison to full MLE was performed. The
theoretical claim applies to $d = 4$; our portfolio has exactly 4 assets.

\asmfield{Consequence if assumption fails}
Slightly biased $\hat{R}$; impact on VaR is negligible ($< 1\%$) for $d = 4$.

\asmfield{Mitigation in this project}
Theoretical guarantee for $d \le 4$; no corrective action needed.

% - - - - - - - - - - - - - - - - - - - - - - - - - - - - - - - - - - - - -
\subsubsection{C.5 --- Cram\'er-von Mises parametric bootstrap is a valid GoF test}\label{asm:C5}

\asmfield{Statement}
The parametric bootstrap CvM test --- comparing the observed CvM statistic
to its null distribution generated by drawing from the fitted $t$-copula ---
is a valid test of whether the data are consistent with the specified
$t$-copula model at the fitted $R$.

\asmfield{Source in code}
\code{var\_copula.py:317-353}: \code{cvm\_gof\_test()} with
$N_{\mathrm{bootstrap}} = 1{,}000$ and vectorised empirical copula evaluation.

\asmfield{Theoretical basis}
\cite{McNeilFreyEmbrechts2005}, \S~7.3.4 (parametric bootstrap GoF).
The test is asymptotically valid under $H_0$ but has low power against
certain departure types (e.g., asymmetric dependence).

\asmfield{Validation in this project}
\diagnostic\ --- Results in \code{outputs/tables/copula\_gof.csv};
rejection at 4 of 18 refits (22.2\%). Test implementation verified correct
via vectorised broadcasting (\code{var\_copula.py:298-314}), avoiding a
Python loop over 500 evaluation points.

\asmfield{Consequence if assumption fails}
Test has low power against asymmetric or vine copula alternatives. Failure
to reject does not guarantee the $t$-copula is correct.

\asmfield{Mitigation in this project}
1{,}000 bootstrap iterations per refit; GoF $p$-values stored per date
in output CSVs.

% - - - - - - - - - - - - - - - - - - - - - - - - - - - - - - - - - - - - -
\subsubsection{C.6 --- Log-return re-pricing formula is correct}\label{asm:C6}

\asmfield{Statement}
The portfolio P\&L for a simulated return vector $\mathbf{r}$ is
$\Delta V_{\text{linear}} = (V_0/d) \sum_{j=1}^{d}(\mathrm{e}^{r_j} - 1)$,
matching the exact arithmetic return from a log-return scenario.

\asmfield{Source in code}
\code{var\_copula.py:429}:
\begin{lstlisting}
PnL_linear = (V0 / N_ASSETS) * (np.exp(r_sim) - 1).sum(axis=1)
\end{lstlisting}

\asmfield{Theoretical basis}
If $\log(P_t/P_{t-1}) = r_j$, then $\Delta P / P = \mathrm{e}^{r_j} - 1$.
For small $r_j$ this approximates $r_j$; for large moves the exact formula
is materially different.

\asmfield{Validation in this project}
\validated\ --- Formula matches \code{VaR\_Estimation.R}, \S~4d (professor's
reference implementation). Previous linear formula ($r_{\text{sim}}$ without
$\exp$) overstated tail losses by $5$--$10\%$ at extreme quantiles;
corrected in commit A1, which reduced copula exceptions from 156 to 51.

\asmfield{Consequence if assumption fails}
Previous formula: 156 exceptions (exception rate 3.65\%, expected 1\%) ---
a factor-of-3.5 over-exception. After fix: 51 exceptions (1.19\%), Kupiec
$p = 0.2149$ (not rejected). The fix is now in production.

\asmfield{Mitigation in this project}
Fix applied and verified against professor's R code (commit A1).

% - - - - - - - - - - - - - - - - - - - - - - - - - - - - - - - - - - - - -
\subsubsection{C.7 --- Nonlinear P\&L is dependence-independent from linear factors}\label{asm:C7}

\asmfield{Statement}
IRS and straddle P\&L can be added to simulated linear P\&L by sampling
an independent historical observation (bootstrap), because the nonlinear
components are not jointly modelled in the copula.

\asmfield{Source in code}
\code{var\_copula.py:432-434}:
\begin{lstlisting}
idx     = rng.integers(0, len(nonlinear_window), size=n_sims)
PnL_sim = PnL_linear + nonlinear_window[idx]
\end{lstlisting}

\asmfield{Theoretical basis}
No theoretical basis --- pragmatic simplification. Full treatment would
add DV01 (IRS rate sensitivity) and delta/gamma/vega (straddle) as
additional copula dimensions.

\asmfield{Validation in this project}
\notvalidated\ --- The code comment at \code{var\_copula.py:54-56} explicitly
flags this as a simplification. No test of nonlinear--linear dependence
was performed.

\asmfield{Consequence if assumption fails}
If IRS losses and equity losses co-move during crises (e.g., flight-to-quality
pushes rates down while equities fall): the independent bootstrap
\emph{underestimates} joint tail losses. The straddle position is long
gamma, which partially offsets this via positive co-movement with equity
volatility. Magnitude: $5$--$15\%$ VaR underestimation in joint stress scenarios.

\asmfield{Mitigation in this project}
Historical bootstrap preserves the empirical mean nonlinear P\&L;
\code{pnl\_total} in the backtest includes the true dependence. The
exception count (51 vs.\ 42.7 expected) suggests mild remaining
undercoverage, consistent with this assumption failing partially.

% - - - - - - - - - - - - - - - - - - - - - - - - - - - - - - - - - - - - -
\subsubsection{C.8 --- $N_{\mathrm{SIMS}} = 20{,}000$ paths give MC SE $\le 0.5\%$ on 99\% VaR}\label{asm:C8}

\asmfield{Statement}
Drawing 20{,}000 MC paths per day is sufficient to estimate the 1st
percentile of the simulated P\&L with a Monte Carlo standard error not
exceeding $0.5\%$ of the VaR estimate.

\asmfield{Source in code}
\code{var\_copula.py:92}: \lstinline|N_SIMS = 20_000|.

\asmfield{Theoretical basis}
MC SE on a quantile $q_\alpha$ with $N$ paths:
$\text{SE}(q_\alpha) \approx \sqrt{\alpha(1-\alpha)} / (N \cdot f(q_\alpha))$,
where $f$ is the PDF at the quantile. For $\alpha = 0.01$ and $N = 20{,}000$,
the SE on the exception rate is $\approx 0.15$ pp --- negligible relative
to the 1\% target.

\asmfield{Validation in this project}
\diagnostic\ --- 3{,}547 distinct daily VaR values across 4{,}269 backtest
days confirm that daily fresh draws produce meaningfully different estimates
(not constant within a refit window). Bootstrap CI for VaR stored in
\code{outputs/tables/var\_bootstrap\_ci.csv}.

\asmfield{Consequence if assumption fails}
Too few paths: VaR estimates are noisy; day-to-day variation is dominated
by MC error rather than genuine model signal. At $N = 5{,}000$ paths the
SE would be $\approx 0.3$ pp on the exception rate --- still acceptable but
noisier.

\asmfield{Mitigation in this project}
Bootstrap CI per refit window provides empirical MC uncertainty quantification.

% - - - - - - - - - - - - - - - - - - - - - - - - - - - - - - - - - - - - -
\subsubsection{C.9 --- Copula refit every 250 days is sufficient}\label{asm:C9}

\asmfield{Statement}
Re-estimating marginals, pseudo-observations, and the copula correlation
matrix every 250 trading days (approximately one year) captures
time-variation in dependence structure adequately; the correlation structure
is stable within each 250-day window.

\asmfield{Source in code}
\code{var\_copula.py:91}: \lstinline|COPULA_REFIT = 250|.

\asmfield{Theoretical basis}
No formal theoretical guidance. Longer refit intervals reduce estimation
noise; shorter intervals increase responsiveness to regime changes.

\asmfield{Validation in this project}
\notvalidated\ --- No sensitivity comparison vs.\ 125-day or 500-day refit
intervals. CvM rejection at 4 of 18 refits suggests the model does not fit
all 250-day windows, which is consistent with correlation instability
within windows.

\asmfield{Consequence if assumption fails}
Major correlation shift (e.g., COVID spike in cross-asset correlations in
March 2020) mid-window: stale $R$ produces wrong copula samples; VaR
under- or overestimated by $\approx 10\%$.

\asmfield{Mitigation in this project}
Rolling 500-day window limits staleness. Sensitivity test listed in
\cref{sec:sensitivity-roadmap}.

% - - - - - - - - - - - - - - - - - - - - - - - - - - - - - - - - - - - - -
\subsubsection{C.10 --- Returns are i.i.d.\ (no GARCH filtering --- known limitation)}\label{asm:C10}

\asmfield{Statement}
Log-returns are treated as i.i.d.\ draws from a stationary distribution
for the purpose of marginal fitting and copula estimation. No GARCH
pre-filtering is applied.

\asmfield{Source in code}
\code{var\_copula.py:38-40}: \emph{``No GARCH filtering --- marginals are
fitted on raw log-returns. Raw returns exhibit volatility clustering which
violates the i.i.d.\ assumption of pseudo-observations; this is a known
limitation of the course approach.''}

\asmfield{Theoretical basis}
\cite{McNeilFrey2000}: volatility clustering is the primary motivation for
GARCH+EVT over unconditional EVT/copula methods.

\asmfield{Validation in this project}
\notvalidated\ --- Explicitly documented as a known limitation. Copula
Christoffersen independence test: $p = 0.0003$ (strongly rejected) ---
exceptions cluster, confirming that the i.i.d.\ assumption fails.

\asmfield{Consequence if assumption fails}
During volatility regime changes: pseudo-observations are not uniform (high-vol
days cluster), biasing $\hat{R}$ upward; copula VaR underestimates tail
risk by an estimated $30$--$50\%$ during peak crisis periods. The 51
exceptions (vs.\ 42.7 expected) and rejected Christoffersen test directly
reflect this failure.

\asmfield{Mitigation in this project}
None within the copula model. This assumption cannot be relaxed without
replacing the entire copula framework with a DCC-GARCH-copula model.
Flagged explicitly in the code docstring.

% ---------------------------------------------------------------------------
\subsection{Category D --- Parameter and Model Choices}\label{sec:catD}
% ---------------------------------------------------------------------------

\subsubsection{D.1 --- Log-returns are the appropriate return definition}\label{asm:D1}

\asmfield{Statement}
Daily portfolio returns are defined as log-returns $r_t = \log(P_t/P_{t-1})$,
and portfolio P\&L is $\Delta V_t \approx V_0 \sum_i w_i r_{i,t}$ (Irle
Eq.~10).

\asmfield{Source in code}
\code{src/data/compute\_returns.py:23}:
\begin{lstlisting}
log_returns = np.log(prices / prices.shift(1)).dropna()
\end{lstlisting}

\asmfield{Theoretical basis}
Irle, p.~35. Log-returns are additive over time; simple returns are additive
across assets. For daily horizons the two are numerically close (difference
$\approx r^2/2 < 0.01\%$ for $|r| < 3\%$).

\asmfield{Validation in this project}
\diagnostic\ --- Standard academic and regulatory practice. No comparison
vs.\ arithmetic returns performed.

\asmfield{Consequence if assumption fails}
For days with $|r| > 5\%$ (crisis days), the approximation error is $> 0.1\%$
per asset per day. For a joint 5\% loss across all assets,
$V_0 (\mathrm{e}^{-0.05} - 1) \approx -\$48{,}770$ vs.\ $-\$50{,}000$
(simple return) --- a $2.5\%$ underestimation of loss.

\asmfield{Mitigation in this project}
Copula model uses exact $(\mathrm{e}^r - 1)$ formula (C.6); linear models
use the log-return approximation (consistent with Irle).

% - - - - - - - - - - - - - - - - - - - - - - - - - - - - - - - - - - - - -
\subsubsection{D.2 --- Equal weights $w_i = 1/N$ across linear assets}\label{asm:D2}

\asmfield{Statement}
Each of the four linear assets (EURUSD, GLD, IEF, SPY) receives equal
weight $w_i = 1/4 = 25\%$ of the portfolio notional.

\asmfield{Source in code}
\code{src/data/config.py:15-20}: \lstinline|WEIGHTS_DICT = {"SPY": 1/4, ...}|.

\asmfield{Theoretical basis}
Portfolio specification (not a statistical assumption). Equal weighting
is a valid baseline; deviations are portfolio management choices, not
statistical errors.

\asmfield{Validation in this project}
\notvalidated\ as a statistical assumption --- this is a portfolio
design choice, not a modeling assumption subject to testing.

\asmfield{Consequence if assumption fails}
Alternative weighting (e.g., risk parity) would change both the P\&L
series and the VaR. Direction and magnitude depend on target weights.

\asmfield{Mitigation in this project}
Weights are hardcoded in \code{config.py} and consistent across all modules.

% - - - - - - - - - - - - - - - - - - - - - - - - - - - - - - - - - - - - -
\subsubsection{D.3 --- $V_0 = \$1{,}000{,}000$ is the relevant scale}\label{asm:D3}

\asmfield{Statement}
The portfolio notional is exactly \$1{,}000{,}000 USD at inception and is
not updated (no mark-to-market rebalancing of the notional).

\asmfield{Source in code}
\code{src/data/config.py:27}: \lstinline|V0 = 1_000_000|.

\asmfield{Theoretical basis}
Standard regulatory VaR reporting convention.

\asmfield{Validation in this project}
\validated\ --- Consistent across all VaR modules by construction.

\asmfield{Consequence if assumption fails}
Wrong $V_0$ scales all VaR numbers proportionally. Not a statistical error.

\asmfield{Mitigation in this project}
Centralised in \code{config.py}; consistent across modules.

% - - - - - - - - - - - - - - - - - - - - - - - - - - - - - - - - - - - - -
\subsubsection{D.4 --- \texttt{pnl\_total} covers all material P\&L sources}\label{asm:D4}

\asmfield{Statement}
\code{data/processed/total\_portfolio\_pnl.csv} correctly aggregates all
material P\&L sources: $\Delta V_t = \Delta V_t^{\text{linear}} +
\Delta V_t^{\text{IRS}} + \Delta V_t^{\text{straddle}}$.

\asmfield{Source in code}
\code{src/data/compute\_pnl.py:79}: \lstinline|total["pnl_total"] = total.sum(axis=1)|.

\asmfield{Theoretical basis}
Complete portfolio P\&L decomposition --- see \code{docs/pricing/00\_overview.md}.

\asmfield{Validation in this project}
\diagnostic\ --- All three components explicitly computed and stored. No
external positions exist outside \code{config.py}.

\asmfield{Consequence if assumption fails}
An unmodelled position would cause VaR to understate true risk.
Magnitude: proportional to the unmodelled position's tail contribution.

\asmfield{Mitigation in this project}
Portfolio fully specified in \code{config.py}; no ad-hoc positions.

% - - - - - - - - - - - - - - - - - - - - - - - - - - - - - - - - - - - - -
\subsubsection{D.5 --- Historical sample period is representative of future regimes}\label{asm:D5}

\asmfield{Statement}
The 2006-01-01 to 2024-12-31 sample contains sufficient market regime
variety to produce VaR estimates applicable to the future distribution of
portfolio losses.

\asmfield{Source in code}
\code{src/data/config.py:23-24}: \lstinline|START_DATE = "2006-01-01"|,
\lstinline|END_DATE = "2024-12-31"|.

\asmfield{Theoretical basis}
Fundamental limitation of all historical-data-based VaR methods. No
theoretical basis for representativeness.

\asmfield{Validation in this project}
\notvalidated\ --- No stress test against scenarios outside the historical
sample (hyperinflation, sovereign default cascade, USD reserve crisis)
was performed.

\asmfield{Consequence if assumption fails}
Tail events outside the 2006--2024 experience would produce losses larger
than the historical maximum ($-\$52{,}940$). All three VaR models would
simultaneously underestimate VaR. Direction: VaR \emph{underestimated};
magnitude: \emph{unknown} by construction.

\asmfield{Mitigation in this project}
None. This is a fundamental limitation of the back-testing methodology.
Regulatory practice typically supplements backtesting with stress scenarios.

% - - - - - - - - - - - - - - - - - - - - - - - - - - - - - - - - - - - - -
\subsubsection{D.6 --- Trading days aligned across all assets}\label{asm:D6}

\asmfield{Statement}
Daily prices for SPY, IEF, GLD, EURUSD=X, VIX, and DGS10 are all
observed on the same set of trading days; no stale-price problem arises
from asynchronous markets.

\asmfield{Source in code}
\code{src/data/compute\_pnl.py:29}:
\begin{lstlisting}
common = spy.index.intersection(vix.index).intersection(dgs10.index)
\end{lstlisting}

\asmfield{Theoretical basis}
Data alignment is a data-quality requirement, not a statistical assumption.
All four equity/FX instruments trade on US market days; DGS10 is a daily
Federal Reserve publication; VIX is CBOE.

\asmfield{Validation in this project}
\validated\ --- Date-intersection logic explicitly takes the common set;
missing dates are excluded. Total aligned days: 4{,}769.

\asmfield{Consequence if assumption fails}
Stale prices on US holidays (EURUSD may trade) introduce spurious zero
returns; VaR is slightly underestimated. Magnitude: negligible (few days
per year).

\asmfield{Mitigation in this project}
Inner-join on all date indices in \code{compute\_pnl.py:29}.

% - - - - - - - - - - - - - - - - - - - - - - - - - - - - - - - - - - - - -
\subsubsection{D.7 --- Pricing assumptions from \texttt{docs/pricing/*.md}}\label{asm:D7}

\asmfield{Statement}
All pricing assumptions documented in \code{docs/pricing/05\_IRS.md} (IRS
DV01 approximation) and \code{docs/pricing/06\_straddle.md} (Black-Scholes
straddle) are accepted without formal validation.

\asmfield{Source in code}
\code{src/data/portfolio\_pricing.py:78-105} (IRS);
\code{src/data/portfolio\_pricing.py:21-76} (straddle).

\asmfield{Theoretical basis}
See \code{docs/pricing/05\_IRS.md} and \code{docs/pricing/06\_straddle.md}
for per-instrument theoretical references.

\asmfield{Validation in this project}
\notvalidated\ formally. IRS DV01 approximation ignores curve shape;
error $\approx 5$--$10\%$ for large parallel shifts. BS straddle ignores
volatility smile/skew; error can exceed $20\%$ for deep in-the-money options,
but ATM straddles are less affected ($\approx 5$--$10\%$).

\asmfield{Consequence if assumption fails}
Mis-priced IRS/straddle produces incorrect \code{pnl\_irs} and
\code{pnl\_straddle} series. All VaR models trained on incorrect P\&L data
will have correspondingly biased estimates.

\asmfield{Mitigation in this project}
Pricing formulas are standard industry approximations. Full curve-based IRS
valuation and smile-adjusted option pricing are out of scope for this project.

% ---------------------------------------------------------------------------
\subsection{Category E --- Backtesting Framework}\label{sec:catE}
% ---------------------------------------------------------------------------

\subsubsection{E.1 --- Kupiec: exception count $\sim \mathrm{Binomial}(T, 1-\alpha)$ under $H_0$}\label{asm:E1}

\asmfield{Statement}
Under the null hypothesis that the VaR model is correctly specified,
the number of VaR exceptions $N$ follows a binomial distribution
$N \sim \mathrm{Bin}(T,\, p)$ where $p = 1 - \alpha$.

\asmfield{Source in code}
\code{backtesting/backtest.py:200-215}: \code{kupiec\_test()} computes
LR$_{\mathrm{UC}} = -2\ln[p^N(1-p)^{T-N} / \hat{p}^N(1-\hat{p})^{T-N}] \sim \chi^2(1)$.

\asmfield{Theoretical basis}
\cite{Kupiec1995}. Irle, Ch.~8, p.~184.

\asmfield{Validation in this project}
\validated\ by construction: the binomial assumption is the definition of
the null hypothesis, not an empirical claim about the data.

\asmfield{Consequence if assumption fails}
The test loses validity if exceptions are serially correlated (addressed by
Christoffersen) or if $T$ is too small to justify the asymptotic $\chi^2$
approximation.

\asmfield{Mitigation in this project}
Christoffersen test (\cref{asm:E2}) complements Kupiec by testing for
serial independence; together they form the standard regulatory backtest
battery.

% - - - - - - - - - - - - - - - - - - - - - - - - - - - - - - - - - - - - -
\subsubsection{E.2 --- Christoffersen: exceptions are independent over time}\label{asm:E2}

\asmfield{Statement}
The Christoffersen LR$_{\mathrm{IND}}$ test assumes that under $H_0$,
consecutive exception indicators $I_t$ are independent (no first-order
Markov dependence).

\asmfield{Source in code}
\code{backtesting/backtest.py:248-292}: \code{christoffersen\_test()},
first-order Markov transition matrix.

\asmfield{Theoretical basis}
\cite{Christoffersen1998}. Irle, Ch.~8.

\asmfield{Validation in this project}
\validated\ by construction (test definition). Empirically: EVT
($p_{\mathrm{IND}} = 0.0064$, rejected) and Copula ($p_{\mathrm{IND}} = 0.0003$,
strongly rejected) --- both fail the independence assumption, confirming
that clustering is present. GARCH+EVT passes ($p_{\mathrm{IND}} = 0.4132$).

\asmfield{Consequence if assumption fails}
If exceptions cluster, the model is not adaptive to changing market conditions;
Kupiec alone would miss this. The Christoffersen test correctly flags EVT
and Copula models.

\asmfield{Mitigation in this project}
Christoffersen CC test ($= $ Kupiec $+$ Independence) is the primary
supervisory test. GARCH pre-filtering addresses the root cause.

% - - - - - - - - - - - - - - - - - - - - - - - - - - - - - - - - - - - - -
\subsubsection{E.3 --- $T = 4{,}269$ days gives sufficient statistical power}\label{asm:E3}

\asmfield{Statement}
The backtest sample of $T = 4{,}269$ trading days provides adequate
statistical power to detect material VaR mis-calibration at the 99\%
confidence level.

\asmfield{Source in code}
\code{backtesting/backtest.py:348}: $T$ determined by the length of the
aligned \code{pnl} and \code{var} series.

\asmfield{Theoretical basis}
Irle, Ch.~8. At $T = 4{,}269$ and $p = 0.01$, the expected exception count
is $42.7$ under $H_0$. Power to detect 50\% over-coverage (exception rate
$1.5\%$ vs.\ target $1\%$): $\approx 80\%$ based on normal approximation
to the binomial.

\asmfield{Validation in this project}
\validated\ by calculation: EVT produces 56 exceptions ($p_{\mathrm{UC}}
= 0.0507$) --- borderline not rejected at $5\%$ despite $31\%$ excess
exceptions. This borderline result is itself evidence that $T = 4{,}269$
is adequate but not generous for detecting moderate mis-calibration.

\asmfield{Consequence if assumption fails}
Shorter backtest ($T < 1{,}000$): power to detect 50\% over-coverage falls
to $\approx 30\%$; Kupiec test has very low sensitivity.

\asmfield{Mitigation in this project}
$T = 4{,}269$ is the maximum available given the 2006--2024 sample and
500-day burn-in window.

% - - - - - - - - - - - - - - - - - - - - - - - - - - - - - - - - - - - - -
\subsubsection{E.4 --- Backtesting against \texttt{pnl\_total} is the correct choice}\label{asm:E4}

\asmfield{Statement}
VaR is a risk measure for the full portfolio P\&L; backtesting against
\code{pnl\_total} (linear $+$ IRS $+$ straddle) is the supervisory-correct
choice rather than against per-instrument or linear-only P\&L.

\asmfield{Source in code}
\code{evt.py:344-356}, \code{garch\_evt.py:541-553}, \code{var\_copula.py}:
all call \code{run\_backtest(pnl=pnl, var=...)} where \code{pnl = pnl\_total}.

\asmfield{Theoretical basis}
Irle, Ch.~8, p.~184: exception indicator $I_t = \mathbf{1}\{P\&L_t <
-\text{VaR}_t\}$ where P\&L is the \emph{full portfolio} daily P\&L.

\asmfield{Validation in this project}
\validated\ --- All three models backtest against the same \code{pnl\_total}
series, ensuring comparability and correctness.

\asmfield{Consequence if assumption fails}
Backtesting against linear-only P\&L would understate exceptions for models
that include nonlinear components, producing artificially good Kupiec
$p$-values.

\asmfield{Mitigation in this project}
Centralized \code{run\_backtest()} function enforces consistent backtesting
across all models.

% ---------------------------------------------------------------------------
\subsection{Category F --- Numerical Implementation}\label{sec:catF}
% ---------------------------------------------------------------------------

\subsubsection{F.1 --- \texttt{scipy genpareto.fit} MLE converges to true parameters}\label{asm:F1}

\asmfield{Statement}
The maximum-likelihood estimator for the GPD shape ($\xi$) and scale
($\sigma$) parameters implemented in \code{scipy.stats.genpareto.fit}
with \code{floc=0} converges to the true parameters in the production
datasets used here.

\asmfield{Source in code}
\code{evt.py:232}: \lstinline|c, loc, sigma = genpareto.fit(exceedances, floc=0)|.

\asmfield{Theoretical basis}
MLE theory: consistency and asymptotic normality hold for $\xi > -0.5$.
Since $\hat{\xi} > 0$ throughout (mean $= 0.0475$), this condition is
comfortably satisfied.

\asmfield{Validation in this project}
\validated\ by absence of failures: zero GPD fit failures recorded in the
4{,}269-day run. Fallback via \code{warnings.warn} (\code{evt.py:275-279})
was never triggered.

\asmfield{Consequence if assumption fails}
Non-converged MLE produces incorrect $\hat{\xi}$ and $\hat{\sigma}$; VaR
is biased in an unknown direction. Empirical fallback mitigates if failure
is detected.

\asmfield{Mitigation in this project}
Try/except block with \code{warnings.warn} and empirical quantile fallback.

% - - - - - - - - - - - - - - - - - - - - - - - - - - - - - - - - - - - - -
\subsubsection{F.2 --- arch library scaling preserves estimation accuracy}\label{asm:F2}

\asmfield{Statement}
Passing returns as percentages (\code{r\_p $\times$ 100}) to \code{arch\_model}
and converting back (\code{cond\_vol / 100}, \code{forecast.variance / 10{,}000})
preserves numerical accuracy of GARCH MLE.

\asmfield{Source in code}
\code{garch\_evt.py:211}: \lstinline|train = r_p.iloc[:t] * 100.0|;
\code{garch\_evt.py:219}: \lstinline|mu_hat = res.params["mu"] / 100.0|;
\code{garch\_evt.py:226}: \lstinline|sigma_sq_curr = fcast.variance.iloc[-1, 0] / 10000.0|.

\asmfield{Theoretical basis}
\code{arch} library documentation: numerical stability of GARCH MLE
requires variance parameters to be $O(1)$; decimal returns ($O(10^{-3})$)
can cause near-zero $\hat{\omega}$ estimates.

\asmfield{Validation in this project}
\validated\ --- Residual mean $= -0.0176$ (expected $\approx 0$), residual
std $= 0.9936$ (expected $\approx 1$), confirming correct unit transformation.

\asmfield{Consequence if assumption fails}
Incorrect scaling would produce biased $\hat{\sigma}_t$ (e.g., 100$\times$
too large), causing VaR to be 100$\times$ overstated. Never occurred in
production.

\asmfield{Mitigation in this project}
Verified by residual diagnostics; division factors hard-coded and reviewed.

% - - - - - - - - - - - - - - - - - - - - - - - - - - - - - - - - - - - - -
\subsubsection{F.3 --- Random seed = 42 produces reproducible MC results}\label{asm:F3}

\asmfield{Statement}
Setting \code{numpy.random.default\_rng(42)} ensures that all Monte Carlo
draws in the copula model are reproducible across runs, and that the chosen
seed does not produce a pathological sample that systematically over- or
underestimates VaR.

\asmfield{Source in code}
\code{var\_copula.py:96}: \lstinline|RANDOM_SEED = 42|;
\code{var\_copula.py} (main): \lstinline|rng = np.random.default_rng(RANDOM_SEED)|.

\asmfield{Theoretical basis}
Law of large numbers: for $N_{\mathrm{SIMS}} = 20{,}000$, any fixed seed
produces a sample that approximates the true distribution to within MC
standard error $\approx 0.15$ pp on the exception rate.

\asmfield{Validation in this project}
\diagnostic\ --- 3{,}547 distinct VaR values across 4{,}269 days (vs.\
$\approx 18$ before commit A2) confirms that fresh daily draws produce
meaningful variation. Full seed-sensitivity test not performed.

\asmfield{Consequence if assumption fails}
Pathological seed: VaR estimate for a specific day deviates by
$> 3 \times \text{MC SE} = 0.45$ pp from the true quantile --- this is a
$(0.3\%)$ probability event per day, negligible.

\asmfield{Mitigation in this project}
$N_{\mathrm{SIMS}} = 20{,}000$ ensures the MC error is small regardless
of seed. Bootstrap CI in \code{outputs/tables/var\_bootstrap\_ci.csv}
quantifies remaining uncertainty.

% - - - - - - - - - - - - - - - - - - - - - - - - - - - - - - - - - - - - -
\subsubsection{F.4 --- Cholesky of $R$ is numerically stable}\label{asm:F4}

\asmfield{Statement}
The sample correlation matrix $R$ of $t_\nu$-transformed pseudo-observations
is positive definite, and Cholesky decomposition succeeds without numerical
instability.

\asmfield{Source in code}
\code{var\_copula.py:283-286}:
\begin{lstlisting}
try:
    L = np.linalg.cholesky(R)
except np.linalg.LinAlgError:
    R = _nearest_pd(R)
    L = np.linalg.cholesky(R)
\end{lstlisting}
\code{\_nearest\_pd()} at \code{var\_copula.py:245-250} implements the
Higham (1988) projection (\cite{Higham1988}).

\asmfield{Theoretical basis}
For well-separated assets with sufficient data ($n = 500 \gg d = 4$), the
sample correlation matrix is almost surely positive definite. Near-singularity
can arise when two assets have near-perfect correlation.

\asmfield{Validation in this project}
\validated\ --- Nearest-PD projection triggered by a warning at
\code{var\_copula.py:267-271} when minimum eigenvalue $< 10^{-8}$; no such
warnings were emitted in production.

\asmfield{Consequence if assumption fails}
Non-PD $R$: Cholesky fails; copula sampling is impossible. Handled by
nearest-PD fallback. The fallback introduces a small perturbation to $R$
but preserves the overall dependence structure.

\asmfield{Mitigation in this project}
Try/except with \code{\_nearest\_pd()} (Higham 1988).

% ===========================================================================
\section{Consequence Matrix}\label{sec:consequence-matrix}
% ===========================================================================

The following table summarises all 40 assumptions sorted by severity of
failure (most consequential first). Columns: \textbf{\#} --- assumption
label; \textbf{Validated?} --- one of \emph{Validated}, \emph{Diagnostic},
or \emph{NOT validated}; \textbf{VaR if violated} --- direction of bias;
\textbf{Magnitude} --- estimated quantitative impact; \textbf{Mitigation}.

\begin{footnotesize}
\begin{longtable}{@{}p{0.65cm} p{4.0cm} p{2.6cm} p{3.4cm} p{2.0cm} p{2.8cm}@{}}
\caption{Consequence matrix --- sorted by severity of failure.}\label{tab:consequence}\\
\toprule
\textbf{\#} & \textbf{Assumption (short)} & \textbf{Validated?} & \textbf{VaR if violated\ldots} & \textbf{Magnitude} & \textbf{Mitigation} \\
\midrule
\endfirsthead
\multicolumn{6}{c}{\tablename\ \thetable\ (continued)} \\
\toprule
\textbf{\#} & \textbf{Assumption (short)} & \textbf{Validated?} & \textbf{VaR if violated\ldots} & \textbf{Magnitude} & \textbf{Mitigation} \\
\midrule
\endhead
\midrule \multicolumn{6}{r}{\emph{continued on next page}} \\
\endfoot
\bottomrule
\endlastfoot

\cref{asm:C10} & Returns i.i.d.\ in copula (no GARCH) & NOT validated & Underestimated during vol.\ clustering & 30--50\% in crises & None; known limitation \\[4pt]

\cref{asm:A4}  & Losses i.i.d.\ in 500d rolling window & NOT validated & Underestimated during regime shifts & 20--50\% in crises & Rolling window partially mitigates \\[4pt]

\cref{asm:D5}  & Historical sample representative & NOT validated & Either direction & Unknown & None; fundamental limitation \\[4pt]

\cref{asm:A1}  & GPD convergence above $u$ & Diagnostic (MEP, Hill) & Underestimated for fat-tail departures & 5--30\% at 99.9\% & Reg.\ floor; $\xi$ clamped \\[4pt]

\cref{asm:C3}  & $t$-copula $\nu=4$ fixed & Diagnostic (CvM, 22\% rejection) & Over- or underestimated & 10--20\% & CvM GoF documented \\[4pt]

\cref{asm:C7}  & Nonlinear P\&L independent of linear & NOT validated & Underestimated in joint stress & 5--15\% & Historical bootstrap used \\[4pt]

\cref{asm:B3}  & GARCH stationarity $\alpha+\beta<1$ & Validated (range 0.97--0.997) & Explosive VaR if violated & Infinite if unit-root & arch enforces constraint \\[4pt]

\cref{asm:B1}  & GARCH(1,1) sufficient & Diagnostic (residual ACF) & Either direction & $<10\%$ & Kupiec $p=0.79$ suggests adequate \\[4pt]

\cref{asm:A2}  & 90th pct is valid GPD threshold & Diagnostic (MEP) & Under/over-estimated & 5--15\% & $N_u\approx 50$ in bounds \\[4pt]

\cref{asm:B4}  & Residuals $\hat{Z}_t$ approx.\ i.i.d. & Diagnostic (QQ-plot) & Underestimated if clustering remains & 10--20\% & EVT step targets residual tail \\[4pt]

\cref{asm:C5}  & CvM bootstrap valid GoF & Diagnostic & Low power vs.\ some alternatives & Unknown & 1{,}000 bootstrap iter.\ used \\[4pt]

\cref{asm:E1}  & Kupiec binomial under $H_0$ & Validated (by construction) & Test invalid; exceptions not detected & Regulatory risk & Christoffersen complements \\[4pt]

\cref{asm:E2}  & Christoffersen: exceptions independent & Validated (by construction) & Model non-adaptive to regime change & Regulatory risk & GARCH pre-filtering \\[4pt]

\cref{asm:B8}  & EVT on residuals reduces bias & Validated (backtest comparison) & Underestimated without GARCH filter & 15--30\% & Two-step is the mitigation \\[4pt]

\cref{asm:F4}  & Cholesky of $R$ stable & Validated (no warnings in prod.) & Simulation fails & Simulation error & Higham 1988 nearest-PD \\[4pt]

\cref{asm:B5}  & GARCH refit every 50d sufficient & NOT validated & Under- or overestimated & $<5\%$ & Expanding window limits drift \\[4pt]

\cref{asm:A3}  & $\xi\in[-0.5,1.0]$ plausible & Supported by literature & Over/underestimated (conservatively) & Rare & Warnings emitted if triggered \\[4pt]

\cref{asm:C1}  & Marginals fit by 6-candidate set & Diagnostic (AIC, KS) & Underestimated if tail heavier & 5--10\% & 6 candidates including $t$ \\[4pt]

\cref{asm:C2}  & Pseudo-obs uniform and independent & Diagnostic (scatter plot) & Biased $\hat{R}$ if marginal wrong & $<5\%$ & Hazen correction applied \\[4pt]

\cref{asm:C4}  & Sample corr.\ $\approx$ MLE for $d\le4$ & Supported by theory & Slightly biased $\hat{R}$ & $<1\%$ for $d=4$ & Theoretical guarantee \\[4pt]

\cref{asm:C6}  & Log-return re-pricing exact & Validated (vs.\ R code; commit A1) & Overstated (before fix: 156 exc.) & Fix reduced to 51 exc. & Exact $e^r-1$ formula used \\[4pt]

\cref{asm:C8}  & $N=20{,}000$ paths sufficient & Diagnostic (bootstrap CI) & Noisy VaR if $N$ too small & MC SE $\approx 0.15$ pp & Bootstrap CI stored \\[4pt]

\cref{asm:C9}  & Copula refit every 250d sufficient & NOT validated & Under/over if correlation shifts & $\approx 10\%$ & Rolling 500d window limits staleness \\[4pt]

\cref{asm:D1}  & Log-returns appropriate & Diagnostic (standard practice) & Small underestimation on crisis days & $<2.5\%$ per large move & Copula uses exact $e^r-1$ \\[4pt]

\cref{asm:D2}  & Equal weights $w=1/4$ & NOT validated (design choice) & Different VaR under alt.\ weights & Depends on weights & Hardcoded in config.py \\[4pt]

\cref{asm:D3}  & $V_0=\$1{,}000{,}000$ correct & Validated (consistent) & Linear VaR scaling error & Linear if wrong & Centralised in config.py \\[4pt]

\cref{asm:D4}  & \code{pnl\_total} complete & Diagnostic (portfolio specified) & Understated if positions missing & Proportional & config.py fully specifies \\[4pt]

\cref{asm:D6}  & Trading days aligned & Validated (inner join) & Stale-price bias on specific days & Negligible & Inner-join in compute\_pnl.py \\[4pt]

\cref{asm:D7}  & Pricing assumptions valid & NOT validated (IRS, BS) & Mis-priced P\&L contaminates VaR & 5--20\% (BS smile) & Standard approximations \\[4pt]

\cref{asm:B2}  & GARCH innovations Student-$t$ & Diagnostic (QQ-plot) & Underestimated if heavier-tailed & $<5\%$ & EVT on residuals compensates \\[4pt]

\cref{asm:B6}  & GARCH $\hat{\mu}$ vs.\ sample mean & NOT validated (negligible) & VaR bias $<0.1\%$ & Negligible & MLE $\hat{\mu}$ is standard \\[4pt]

\cref{asm:B7}  & $\hat{\sigma}_t$ correct conditioning & Validated (structural) & Follows B.1 & See B.1 & Code review verified \\[4pt]

\cref{asm:E3}  & $T=4{,}269$ sufficient power & Validated (power calc.) & Low power if $T<1{,}000$ & $\approx 80\%$ power at $T=4{,}269$ & Max available sample \\[4pt]

\cref{asm:E4}  & Backtest vs.\ \code{pnl\_total} correct & Validated (all models use same) & Understated exceptions if wrong & Systematic bias & Centralised run\_backtest() \\[4pt]

\cref{asm:A5}  & KS $p$-value informative & NOT validated (known bias) & Over-permissive quality indicator & 20--40\% upward bias & Used as relative indicator only \\[4pt]

\cref{asm:A6}  & Empirical quantile fallback OK & Validated (never triggered) & VaR bounded by sample max & Moot (0 activations) & $N_u$ guard prevents activation \\[4pt]

\cref{asm:A7}  & Regulatory floor defensible & Validated (audit trail) & Conservative overestimation & 0 binding days & \code{VaR\_raw} preserved \\[4pt]

\cref{asm:F1}  & scipy genpareto.fit converges & Validated (0 failures) & Biased $\hat{\xi}$, $\hat{\sigma}$ & Never occurred & Fallback on exception \\[4pt]

\cref{asm:F2}  & arch scaling $\times100/\div10{,}000$ & Validated (residual diagnostics) & VaR $100\times$ wrong if scaling off & Never occurred & Residual std $\approx 1$ verified \\[4pt]

\cref{asm:F3}  & Seed=42 not pathological & Diagnostic (3{,}547 distinct VaR) & VaR off by $\le 3\times$ MC SE & $\le 0.45$ pp & $N=20{,}000$ makes MC SE small \\[4pt]

\end{longtable}
\end{footnotesize}

% ===========================================================================
\section{Sensitivity Analysis Roadmap}\label{sec:sensitivity-roadmap}
% ===========================================================================

The following sensitivity analyses are \emph{not} performed in the current
project but would formally test the listed assumptions. They are presented
to demonstrate methodological awareness and to guide future validation effort.

\medskip
\begin{tabularx}{\textwidth}{@{}p{3.8cm} X p{2.0cm}@{}}
\toprule
\textbf{Assumption tested} & \textbf{Sensitivity test} & \textbf{Effort} \\
\midrule

GARCH(1,1) sufficient (\cref{asm:B1}) &
Re-run \code{garch\_evt.py} with \code{GARCH(2,1)} and \code{GARCH(1,2)}
using \code{arch\_model(..., p=2, q=1)} and \code{(p=1, q=2)};
compare backtest exception counts and Kupiec $p$-values; flag if
GARCH(2,1) produces materially better Christoffersen result &
1--2 h \\[4pt]

Threshold $q = 90\%$ (\cref{asm:A2}) &
Re-run \code{evt.py} and \code{garch\_evt.py} with
\code{THRESHOLD\_Q} $\in \{0.85, 0.925, 0.95\}$;
plot rolling $\xi$ time series for each; confirm GPD fit stability
via KS $p$-value distribution; compare mean VaR and exception counts &
30 min \\[4pt]

Copula $\nu$ fixed at 4 (\cref{asm:C3}) &
Re-run \code{var\_copula.py} with \code{NU\_T\_COPULA}
$\in \{3, 6, 10, 30\}$; report mean VaR, exception count,
Kupiec $p$-value, and tail-dependence coefficient
$\lambda_U = 2\,t_{\nu+1}(-\sqrt{(\nu+1)(1-\rho)/(1+\rho)})$
for each $\nu$ &
1 h \\[4pt]

MC paths sufficient (\cref{asm:C8}) &
Re-run \code{monte\_carlo\_var} with \code{N\_SIMS}
$\in \{5{,}000, 100{,}000\}$; compare daily VaR distributions
and bootstrap CI widths; confirm that $N = 20{,}000$ lies in the
stable region where MC SE is negligible &
30 min \\[4pt]

GARCH refit cadence (\cref{asm:B5}) &
Re-run \code{fit\_garch\_expanding} with \code{REFIT\_EVERY}
$\in \{25, 100, 250\}$; compare parameter time series
($\alpha_1 + \beta$ drift) and backtest results;
identify whether faster refit improves Christoffersen $p$-value &
2--3 h \\[4pt]

Copula refit cadence (\cref{asm:C9}) &
Re-run \code{compute\_copula\_var} with \code{COPULA\_REFIT}
$\in \{125, 500\}$; compare CvM rejection rates and exception counts;
flag whether CvM rejection rate changes with window length &
1--2 h \\[4pt]

Rolling vs.\ expanding window (\cref{asm:A4}, \cref{asm:C10}) &
Re-run \code{evt.py} with \code{WINDOW}
$\in \{250, 750, 1{,}000\}$; compare rolling vs.\ expanding for
tail stability, mean $\xi$, and responsiveness to crises;
most informative sensitivity for EVT specifically &
1 h \\[4pt]

Equal vs.\ optimised weights (\cref{asm:D2}) &
Recompute \code{WEIGHTS\_DICT} using minimum-variance or risk-parity
optimisation; re-run full pipeline; compare mean VaR and exception count &
2 h \\[4pt]

Log vs.\ simple returns (\cref{asm:D1}) &
Replace log-return P\&L approximation with arithmetic P\&L
$\Delta V = V_0 \sum_i w_i (P_{i,t}/P_{i,t-1} - 1)$;
compare daily VaR on the 50 worst-loss days &
1 h \\[4pt]

MC seed sensitivity (\cref{asm:F3}) &
Re-run \code{var\_copula.py} with seeds
$\{0, 1, 2, 100, 99{,}999\}$; compare daily VaR series;
confirm standard deviation of seed-to-seed VaR variation
is within MC SE bounds &
2 h \\[4pt]

\bottomrule
\end{tabularx}

% ===========================================================================
\section{Cross-References}\label{sec:cross-references}
% ===========================================================================

\subsection{Links to pricing documentation}

\cref{asm:D7} cross-references the following pricing assumption documents
(Markdown sources in \code{docs/pricing/}; no \LaTeX\ versions exist):

\begin{itemize}
  \item \code{docs/pricing/00\_overview.md} --- portfolio structure, all 6 instruments
  \item \code{docs/pricing/01\_SPY.md} --- log-return P\&L formula (see \cref{asm:D1})
  \item \code{docs/pricing/02\_IEF.md} --- log-return P\&L (duration risk implicit)
  \item \code{docs/pricing/03\_GLD.md} --- log-return P\&L (commodity)
  \item \code{docs/pricing/04\_EURUSD.md} --- log-return P\&L (FX spot)
  \item \code{docs/pricing/05\_IRS.md} --- DV01 approximation for IRS mark-to-market
  \item \code{docs/pricing/06\_straddle.md} --- Black-Scholes straddle with constant
    risk-free rate $r_f = 5\%$ and VIX as implied volatility
\end{itemize}

Key pricing assumption risks:
\begin{itemize}
  \item IRS DV01 approximation (\code{portfolio\_pricing.py:100}): single-discount-factor
    formula accurate to $\approx 5$--$10\%$ for large parallel shifts; acceptable for
    daily P\&L computation.
  \item ATM straddle Black-Scholes (\code{portfolio\_pricing.py:43-50}): ignores
    volatility smile/skew; for ATM options, BS is accurate to $\approx 5\%$; error
    grows for moneyness deviations as the straddle ages.
  \item Constant risk-free rate (\code{config.py:38}): $r_f = 5\%$ throughout;
    see \cref{asm:E1} (market/economic assumptions). During 2008--2022 the true
    risk-free rate ranged from $0\%$ to $5.5\%$; impact on straddle pricing is
    secondary to SPY and VIX moves.
\end{itemize}

\subsection{Regulatory floor decisions}

The regulatory floor $\max(\text{VaR}, u, 0)$ in both EVT and GARCH+EVT
is documented at:
\begin{itemize}
  \item \code{evt.py:262-264} --- EVT floor (\cref{asm:A7})
  \item \code{garch\_evt.py:431} --- GARCH+EVT residual-quantile floor
  \item \code{garch\_evt.py:504} --- Final VaR floor at 0
  \item \code{var\_copula.py:438} --- Copula VaR floor at 0
\end{itemize}
All floors are applied to the \emph{policy output} (\code{VaR} column);
the raw statistical estimator is preserved in separate \code{VaR\_raw} /
\code{VaR\_EVT\_raw} / \code{VaR\_GARCH\_EVT\_raw} columns.

\subsection{Backtest results cited in this appendix}

\begin{center}
\begin{tabular}{@{}lcccccc@{}}
\toprule
\textbf{Model} & \textbf{$T$} & \textbf{$N_{\text{obs}}$} & \textbf{$N_{\text{exp}}$} &
\textbf{Kupiec $p$} & \textbf{Christo.\ $p_{\text{IND}}$} & \textbf{CC $p$} \\
\midrule
EVT            & 4{,}269 & 56 & 42.7 & 0.0507 & 0.0064 & 0.0036 \\
GARCH+EVT      & 4{,}269 & 41 & 42.7 & 0.7936 & 0.4132 & 0.6914 \\
Copula         & 4{,}269 & 51 & 42.7 & 0.2149 & 0.0003 & 0.0007 \\
\bottomrule
\end{tabular}
\end{center}

Source files: \code{outputs/tables/backtest\_evt.csv},
\code{outputs/tables/backtest\_garch\_evt.csv},
\code{outputs/tables/backtest\_copula.csv}.

% ===========================================================================
\section*{Self-Audit}\label{sec:audit}
\addcontentsline{toc}{section}{Self-Audit}
% ===========================================================================

\textbf{1.\ Total assumptions documented by category:}
A (distributional): 7; B (independence/temporal): 8;
C (stationarity/copula): 10; D (parameter/model): 7;
E (backtesting): 4; F (numerical): 4. Total: \textbf{40}.

\textbf{2.\ Validation status:}
Formally validated: 9 (A.6, A.7, B.3, B.7, B.8, C.6, D.3, D.6, E.4 --- plus
F.1, F.2 = 11 total validated). Diagnostic only: 16. NOT validated: 13.

\textbf{3.\ Code files scanned:} 8 (listed in \cref{sec:scope-discovery}).

\textbf{4.\ Files NOT scanned:} delta\_normal.py, historical\_simulation.py
(other team members); download\_data.py (no modeling assumptions); all
\code{plot\_*} functions (presentational only);
\code{docs/pricing/\_generate\_docx.py} (no model logic).

\textbf{5.\ Consequence Matrix orphan check:} All 40 assumption labels
(A.1--A.7, B.1--B.8, C.1--C.10, D.1--D.7, E.1--E.4, F.1--F.4) appear
in both the per-assumption entries (\cref{sec:entries}) and the Consequence
Matrix (\cref{tab:consequence}). No orphans.

\textbf{6.\ Cross-references to pricing docs:} All 7 files listed in
\cref{sec:cross-references} exist under \code{docs/pricing/*.md}.
No \LaTeX\ versions (\code{*.tex}) exist; cross-references point to Markdown
sources. This gap is documented explicitly.

\textbf{7.\ Compilation status:} \code{pdflatex} not found in this environment.
The \code{.tex} file is complete and standalone. To compile:
\begin{lstlisting}[language=bash]
cd docs/appendix
pdflatex 00_assumption_appendix.tex
bibtex   00_assumption_appendix
pdflatex 00_assumption_appendix.tex
pdflatex 00_assumption_appendix.tex
\end{lstlisting}
Expected unresolved references before BibTeX: 10 \verb|\cite{}| calls.
After BibTeX: all resolved. Known potential warning: \code{longtable}
page-break behaviour with \code{hyperref} --- add \code{\textbackslash usepackage{hypcap}}
if cross-reference anchors shift.

\textbf{8.\ Citation count:} 10 \verb|\cite{}| commands referencing:
Pickands1975, BalkemadeHaan1974, McNeilFreyEmbrechts2005 (3$\times$),
McNeilFrey2000 (2$\times$), Kupiec1995, Christoffersen1998, Bollerslev1986,
Higham1988. All defined in \code{references.bib}.

% ===========================================================================
\bibliographystyle{plain}
\bibliography{references}
% ===========================================================================

\end{document}
